% Chapter 4 - CLARION (the proposed system)
% Design-level description of the artifact. NO code identifiers (no filenames, functions, src/).
% Two diagrams only: a component architecture map and a single connected end-to-end flow. EN-GB.

\chapter{CLARION: An Explainable Financial-Alert System for Retail Investors}
\label{chap:clarion}

This chapter presents \textsc{Clarion}, the system built in this dissertation: its architecture, the
historical knowledge base and correlation engine at its core, the anomaly trigger, the explanation engine
that makes every alert traceable, and the way alerts reach the investor. Chapter~\ref{chap:methods}
argued for the individual methods; the focus here is how they fit together into one working system.

\section{Overview}
\label{sec:clarion_overview}

\textsc{Clarion} watches the US equity market on behalf of a retail investor and raises an alert
whenever one of two independent triggers fires: an abrupt price movement, or a relevant incoming news
item. Each alert is accompanied by a full, traceable explanation rather than a bare notification. The
system is organised as a pipeline of well-separated components, each with a single responsibility,
connected by a common data schema so that the live and historical layers are directly comparable.

\section{Architecture}
\label{sec:clarion_architecture}

The system is driven by the two triggers and converges on a single explanation step before delivery.
Figure~\ref{fig:architecture} shows the components and the flow of data between them. The live layer
supplies real-time prices and news, and the historical knowledge base supplies the precedents. The
anomaly detector and the correlation engine turn these into evidence, the explanation engine assembles a
traceable account, and the delivery component sends the alert to the investor.

\begin{figure}[ht]
\centering
\begin{tikzpicture}[
  font=\small,
  >={Stealth[length=2.2mm]},
  srcbox/.style={draw, rounded corners, align=center, text width=2.9cm, minimum height=1.0cm, fill=black!8},
  comp/.style={draw, rounded corners, align=center, text width=2.9cm, minimum height=1.0cm, fill=black!3},
  node distance=1.15cm and 1.35cm
]
  % data sources (top)
  \node[srcbox] (md) {Market data\\ \footnotesize live prices};
  \node[srcbox, right=of md] (nf) {News feed\\ \footnotesize live news};
  \node[srcbox, right=of nf] (kb) {Historical\\ knowledge base};
  % engines
  \node[comp, below=of md] (an) {Anomaly detector\\ \footnotesize \emph{z}-score};
  \node[comp, below=of nf] (ce) {Correlation engine\\ \footnotesize embeddings + event study};
  % explanation (centred between the two engines)
  \node[comp] (ex) at ($(an)!0.5!(ce) + (0,-1.7)$) {Explanation engine\\ \footnotesize transparent, faithful};
  % delivery
  \node[comp, below=1.15cm of ex] (tg) {Alert delivery\\ \footnotesize Telegram};
  % flows
  \draw[->] (md) -- (an);
  \draw[->] (nf) -- (ce);
  \draw[->] (kb) -- (ce);
  \draw[->] (an) -- (ex);
  \draw[->] (ce) -- (ex);
  \draw[->] (ex) -- (tg);
\end{tikzpicture}
\caption[System architecture and data flow]{Architecture of \textsc{Clarion}. The live layer (market
data, news feed) and the historical knowledge base feed two engines (the anomaly detector and the
correlation engine) whose evidence converges on the explanation engine before an alert is delivered
through Telegram. The market-movement trigger follows the left path; the news trigger the right.}
\label{fig:architecture}
\end{figure}

Three engineering principles recur throughout and make the system both testable and defensible. First,
\textbf{a thin end-to-end slice was built before breadth}: the market-movement trigger was implemented
and proven end to end (from price retrieval to a real alert) before the heavier correlation
engine was added, reducing integration risk. Second, \textbf{pure logic is separated from
input/output}, so the numerical core can be verified deterministically and offline, without a network
connection or the heavy model stack. Third, \textbf{components are programmed against interfaces}: the
embedding step is defined abstractly, with two interchangeable implementations (a dependency-free
lexical baseline and the real Sentence-BERT model), so that one can be substituted for the other
without changing anything else, which is what makes the embedding-model comparison in
Chapter~\ref{chap:casestudies} a fair one.

Table~\ref{tab:clarion_decisions} gathers the principal design decisions and the rationale behind each,
making explicit the engineering judgement that, in a work of this kind, \emph{is} the contribution.
Every decision follows the guiding principle of defensible simplicity established in
Chapter~\ref{chap:methods}.

\begin{table}[ht]
\caption{Principal design decisions in \textsc{Clarion} and their rationale.}
\label{tab:clarion_decisions}
\centering
\begin{tabular}{p{0.27\textwidth} p{0.30\textwidth} p{0.33\textwidth}}
\toprule
\tabhead{Decision} & \tabhead{Choice} & \tabhead{Rationale} \\
\midrule
Detection method & Rolling \emph{z}-score & Transparent; every alert is explainable to a non-expert \\
News representation & Sentence-BERT embeddings & Contextual meaning, yet cheap and reproducible to compare \\
Impact measurement & Event-study windows & Standard, reproducible; characterises an outcome, not a forecast \\
Explanation strategy & Transparent design + precedents & Faithful by construction, not a post-hoc approximation \\
Data architecture & Shared schema, two layers & Live and historical evidence are directly comparable \\
Component coupling & Programmed to interfaces & Testable offline; enables fair model substitution \\
Delivery channel & Telegram bot & Free, ubiquitous, no bespoke client needed \\
\bottomrule
\end{tabular}
\end{table}

Where Figure~\ref{fig:architecture} shows the components, Figure~\ref{fig:workflow} traces the same
system as one connected flow of data and processing steps: the offline knowledge-base build (top), the
news trigger (middle) and the market trigger (bottom), all converging on the explanation engine and the
alert.

\begin{figure}[ht]
\centering
\begin{tikzpicture}[
  font=\scriptsize, >={Stealth[length=2mm]},
  box/.style={draw, rounded corners, align=center, text width=1.7cm, minimum height=0.8cm},
  db/.style={box, fill=black!8},
  pr/.style={box, fill=black!3},
  dec/.style={draw, diamond, aspect=2, align=center, inner sep=1pt, text width=1.0cm, fill=black!5},
  lane/.style={font=\itshape, text=black!55, anchor=east},
]
  % offline build (top)
  \node[db] (o1) at (2.3,0) {FNSPID\\ news $+$ prices};
  \node[db] (kb) at (5.5,0) {Knowledge\\ base};
  \draw[->] (o1) -- node[midway, above, font=\tiny, align=center]{align,\\ embed} (kb);
  % news trigger (middle)
  \node[db] (n1) at (0.9,-1.9) {Live\\ news};
  \node[pr] (n2) at (3.2,-1.9) {SBERT\\ embed};
  \node[pr] (n3) at (5.5,-1.9) {cosine\\ top-$k$};
  \node[pr] (n4) at (7.8,-1.9) {precedents\\ $+$ impact};
  \draw[->] (n1)--(n2); \draw[->] (n2)--(n3); \draw[->] (n3)--(n4);
  \draw[->] (kb) -- node[right, font=\tiny]{precedents} (n3);
  % market trigger (bottom)
  \node[db] (m1) at (0.9,-3.8) {Live\\ prices};
  \node[pr] (m2) at (3.2,-3.8) {log-returns};
  \node[pr] (m3) at (5.5,-3.8) {rolling\\ \emph{z}-score};
  \node[dec] (m4) at (7.8,-3.8) {$|z|{>}k$?};
  \draw[->] (m1)--(m2); \draw[->] (m2)--(m3); \draw[->] (m3)--(m4);
  % converge on explanation + delivery
  \node[pr] (ex) at (10.1,-2.85) {Explanation\\ engine};
  \node[pr] (tg) at (10.1,-4.35) {Telegram\\ alert};
  \draw[->] (n4) -- (ex);
  \draw[->] (m4) -- node[below=1pt, font=\tiny]{yes} (ex);
  \draw[->] (ex) -- (tg);
  % lane labels
  \node[lane] at (-0.2,0)    {offline};
  \node[lane] at (-0.2,-1.9) {news};
  \node[lane] at (-0.2,-3.8) {market};
\end{tikzpicture}
\caption[End-to-end data and processing flow]{The end-to-end flow of \textsc{Clarion} as one connected
pipeline. \emph{Offline} (top), FNSPID news is aligned to prices and embedded into the knowledge base.
The \emph{news trigger} (middle) embeds an incoming headline, finds its nearest precedents in the
knowledge base and attaches their observed impact; the \emph{market trigger} (bottom) turns live prices
into a \emph{z}-score and fires only when it exceeds the threshold. Both paths converge on the
explanation engine, which delivers one explained alert through Telegram.}
\label{fig:workflow}
\end{figure}

\section{Historical Knowledge Base and Correlation Engine}
\label{sec:clarion_kb}

The historical layer is a knowledge base of past news, each entry storing its date, ticker, headline,
sentence embedding and measured post-event impact. Each news item is aligned to the market by taking
the first trading day on or after the news date and measuring impact only from that day's close onward;
this is the no-lookahead rule made concrete, and it avoids crediting the system with a jump already
embedded in the opening price. The intended primary source is the FNSPID corpus
\autocite{dong2024fnspid}; because building the full multi-year corpus is a long one-off job, the
evaluation in Chapter~\ref{chap:casestudies} uses an initial knowledge base built from real recent news
through the identical build process, and the full build is identified as future work.

The correlation engine has two parts. Similarity is computed with cosine over the (normalised)
embeddings, returning the nearest historical items to a query. Impact is computed as the cumulative
return at each horizon ($+1$, $+3$, $+5$ days) measured from the event close, with horizons that fall
outside the available series reported as not available rather than silently dropped. Given a new item,
the engine embeds it, ranks the knowledge base by similarity, and returns the top precedents with their
scores, the evidence later shown to the user. The middle lane of Figure~\ref{fig:workflow} traces these
steps end to end, from an incoming headline to the delivered alert: the headline is embedded, matched
against the knowledge base by similarity, and its nearest precedents are returned with their measured
impact, ready for the explanation engine. The dashed precedents arrow makes plain that the knowledge
base is consulted but never decides; the reasoning stays in the open.

\section{Anomaly Trigger}
\label{sec:clarion_anomaly}

The anomaly trigger applies Equation~\ref{eq:zscore} directly. For the day under evaluation it computes
the mean and standard deviation over the window of strictly preceding days, forms the $z$-score of the
latest return, and returns the decision together with every quantity that produced it: the
$z$-score, window, threshold, mean and standard deviation. Exposing these values is what allows the
explanation engine to be faithful by construction, and the evaluation applies exactly this rule across
entire price series so that the experiments use the same logic as production. The bottom lane of
Figure~\ref{fig:workflow} traces the market trigger from a daily price to the delivered alert: the
latest return is standardised against its rolling window, and the move, with all the quantities behind
it, is passed on only when the absolute $z$-score exceeds the threshold.

\section{Explanation Engine}
\label{sec:clarion_xai}

The explanation engine renders alerts directly from the objects produced upstream, which guarantees
that the text never diverges from the computation. For the market trigger, it states the move, the
$z$-score, the threshold and window, and the recent mean and standard deviation, and it renders the
$z$-score in plain language (how many times a typical day's move it represents), so that a non-expert can
interpret the alert rather than merely see the numbers. For the news trigger,
it lists the retrieved precedents (date, ticker, similarity score, measured impact and headline) and
the average impact across them, and it always ends with an explicit disclaimer that the figures are
observed past outcomes, not a prediction, keeping the system within its no-forecasting scope.

\section{Alert Delivery}
\label{sec:clarion_delivery}

Alerts are delivered to the investor through the Telegram messaging platform, chosen because it offers a
free, widely used bot interface. The two end-to-end flows (retrieve, detect, explain and send for the
market trigger; embed, retrieve precedents, explain and send for the news trigger) each encapsulate a
complete run; periodic scheduling and deployment are deliberately left outside the present scope. A real
alert was successfully delivered during testing; Figure~\ref{fig:telegram_mockup} reproduces the format
the investor receives, with the full reasoning chain visible in a single message.

\begin{figure}[ht]
\centering
\footnotesize
\setlength{\fboxsep}{6pt}%
\colorbox{black!4}{%
  \begin{minipage}{9.2cm}
    \colorbox{black!15}{\makebox[\linewidth][l]{\textbf{CLARION}\, \textemdash{}\, financial alerts bot}}\\[5pt]
    \fcolorbox{black!30}{white}{%
      \begin{minipage}{8.4cm}
        \medskip
        \textbf{News alert \textemdash{} NVDA}\\[2pt]
        \textit{``Nvidia raises guidance on surging demand for AI data-centre accelerators''}\\[4pt]
        \textbf{Potential impact} (across the 3 precedents shown): average 1-day move \textbf{$-2.2\%$}\\[4pt]
        \textbf{Historical precedents:}\\
        $\bullet$~MSFT \ (sim 0.68) \ +1d $-3.5\%$\\
        $\bullet$~META \ (sim 0.68) \ +1d $-2.7\%$\\
        $\bullet$~GOOGL \ (sim 0.64) \ +1d $-0.5\%$\\[4pt]
        {\scriptsize\emph{Precedents are retrieved by semantic similarity; the impact is the observed
        past outcome, not a price prediction.}}
        \medskip
      \end{minipage}}
  \end{minipage}}
\normalsize
\caption[Telegram alert as received by the investor]{A CLARION news alert as delivered through Telegram.
A single message carries the full, traceable reasoning chain: the detected event, the average observed
impact of analogous past events, the individual precedents with their similarity scores and measured
impact, and an explicit no-prediction disclaimer. The figures shown are a representative, rounded
illustration; the full real end-to-end output appears in Chapter~\ref{chap:casestudies} (Case Study~3).}
\label{fig:telegram_mockup}
\end{figure}

\section{Deployment and Practical Use}
\label{sec:clarion_deploy}

CLARION runs as two entry points, one per trigger, each of which performs a complete cycle. The market
trigger fetches recent prices, computes returns, detects, explains and sends; the news trigger takes a
headline, embeds it, retrieves precedents, explains and sends. In day-to-day use these would be invoked
on a schedule (the market trigger once per trading day after the close, the news trigger whenever a
watched feed publishes), but the scheduler and the hosting fall outside the project's scope. The entry
points already encapsulate everything such a scheduler would call.

Running the system needs three things, all free: a market-data source for prices, a company-news or RSS
source for headlines (an API key, kept out of version control), and the historical knowledge base, built
once from the corpus. Appendix~\ref{app:reproducibility} gives the exact steps. The real alert delivered
during testing shows that the path from raw data to the investor's phone works end to end, not just on
paper.

Two practical limits are worth stating plainly. First, the free news tier returns only a recent window,
so the knowledge base is shallower than the multi-year build it is designed for; this affects coverage,
not the method. Second, CLARION is a decision-support prototype rather than a production service: it has
no user accounts, no persistence beyond the knowledge-base file, and no automatic recovery from a
provider outage. None of these is a research limitation; each is a known, ordinary step on the road from
prototype to product, and is recorded here openly rather than hidden.

Beyond these engineering gaps, an alerting tool aimed at non-experts raises two responsible-design
questions that the prototype surfaces rather than solves. The first is alert fatigue: a tool that fires
too often trains the user to ignore it, defeating its purpose. A production version would rank alerts by
severity (for the market trigger, the magnitude of the $z$-score) and rate-limit the least important
notifications, so that attention is reserved for genuinely unusual events. The second is over-reliance: a
confident-looking explanation can be trusted more than its evidence warrants, the failure mode reviewed in
Section~\ref{sec:sota_trust}. The present design mitigates this by showing the individual precedents and
their spread rather than a lone averaged number, and by disclaiming prediction on every alert; natural
extensions are to de-duplicate near-identical precedents and to flag when a retrieved cluster disagrees in
direction, so the user sees genuine uncertainty instead of false consensus. Both are recorded as design
considerations and future work, not as implemented features.

\section{Chapter Conclusions}
\label{sec:clarion_summary}

This chapter described CLARION at the level of its architecture and design: a two-trigger, two-layer
pipeline whose components are separated by responsibility and connected by a common schema, with
explainability built in by rendering every alert directly from the system's own computation. The next
chapter evaluates the system's two core components and its end-to-end behaviour on real data.
